% ============================================================
% M3 成卷轮 S4 件1 —— 单元素养测评卷（二）第二章（19 题足额·新高考 19 题式）
% 2026-09-14｜执行：M3 S4 卷件生产臂
% 骨架＝M2 测评卷 main.tex 件型层 A~E 照抄；卷名/副题/键式改 M3（2章-测-N）。
% 题源：卷面蓝图 §一 19 席矩阵（口袋账 3｜账内 11｜扩容 2｜命制 3）；
%   S5-W4 阶段二重产（0914）：案B 附卷制退役——三栏回流制 \juancols，答案块随题入流
%   （题后紧跟）；卷末附卷页／\anskey 补偿发射／速查表 \dabiao 全撤，锚点 ansblock 恒发。
% 形态转换（生产报告「形态转换登记」节留痕）：测4 填空→单选（双解封装 C）；
%   测12 单选→填空（蓝图授权「两形任取，填空呈现」）；测17 填空→解答；
%   测18 填空→解答；测19 单选→解答（最小改设问，题干逐字）。
% 编译门：xelatex main.tex ×2 / main-pure.tex ×2（sishou remember picture 硬依赖）。
% ============================================================
\documentclass[fontset=none]{ctexart}
\usepackage{qp-m3}
\ansblockgrayfalse   % 换装（测评/滚动＝8 开三栏件）：答案块一律括线模（方案草案§1.2.5/§5.1）
\providecommand{\ov}[1]{\overrightarrow{#1}}

% ------------------------------------------------------------
% 件型层 A｜页面：8 开横放三栏（照 v11 样张件型层，零改动）
% ------------------------------------------------------------
\geometry{paperwidth=399.8mm,paperheight=284.2mm,left=8.9mm,right=8.9mm,
  top=12.0mm,bottom=17.6mm,twoside,headheight=0pt,headsep=0pt,footskip=7.1mm}
\setlength{\topskip}{0pt}
\raggedbottom
\renewcommand{\normalsize}{\fontsize{10.5pt}{17.7pt}\selectfont}
\linespread{1}\normalsize
\setlength{\parindent}{0pt}
\setlength{\parskip}{0pt}
\newlength{\jpcolw}\setlength{\jpcolw}{120.0mm}
\newlength{\jpgap}\setlength{\jpgap}{11.0mm}
\xeCJKsetup{CJKglue={\hskip 0.04em plus 0.04em minus 0.01em}}

\definecolor{silklight}{gray}{0.8235}
\definecolor{silkdark}{gray}{0.6510}
\definecolor{juanrule}{gray}{0.1255}
\definecolor{subgrey}{gray}{0.4}

% ------------------------------------------------------------
% 件型层 B｜JT-01 丝带角饰（照抄样张）
% ------------------------------------------------------------
\newcommand{\juanmathSize}{17pt}
\newcommand{\sishou}{%
  \begin{tikzpicture}[overlay,remember picture,x=1mm,y=1mm]
    \path (current page.north west) coordinate (O);
    \begin{scope}[shift={(O)}]
      \fill[silklight] (6.9,-9.39) -- (40.1,-9.39) -- (6.9,-42.59) -- cycle;
      \draw[white,line width=0.28mm] (9.21,-40.29) -- (9.21,-11.63) -- (37.87,-11.63);
      \fill[silkdark] (22.94,-8.21) -- (34.15,-8.21) -- (5.78,-36.58) -- (5.78,-25.37) -- cycle;
      \fill[silkdark] (34.15,-8.21) -- (34.15,-9.20) -- (33.16,-9.20) -- cycle;
      \fill[silkdark] (5.78,-36.58) -- (6.78,-36.58) -- (6.78,-35.58) -- cycle;
      \node[white,inner sep=0pt,anchor=center,rotate=45] at (17.23,-19.74)
        {\fontsize{\juanmathSize}{\juanmathSize}\selectfont\heitiao 数学};
    \end{scope}
  \end{tikzpicture}}

% ------------------------------------------------------------
% 件型层 C｜卷头零件＋题区零件（照抄样张；卷名（二）/副题第二章）
% ------------------------------------------------------------
\newcommand{\zeroheight}[1]{\raisebox{0pt}[0pt][0pt]{#1}}
\newcommand{\jpcen}[1]{\hbox to\jpcolw{\hfil\zeroheight{#1}\hfil}\nointerlineskip}
\newcommand{\jpright}[1]{\hbox to\jpcolw{\hfill\zeroheight{#1}}\nointerlineskip}
\newcommand{\vgt}[1]{\par\vskip\dimexpr#1-6.25mm\relax}

\newcommand{\tianxieg}[1]{{\fontsize{\qpJTbLabelSize}{13pt}\selectfont\heitiao #1：}%
  \rule[-0.62mm]{\qpJTbLineLen}{0.2mm}}
\newcommand{\jtianxie}{\tianxieg{班级}\hspace{\qpJTbGapGrp}\tianxieg{姓名}%
  \hspace{\qpJTbGapGrp}\tianxieg{得分}}
\newcommand{\juanmingSize}{20.5pt}
\newcommand{\jjuanming}{{\fontsize{\juanmingSize}{26pt}\selectfont\heizhang
  \renewcommand{\CJKglue}{\hskip 0pt}单元素养测评卷（二）}}
\newcommand{\jjunrule}{{\color{juanrule}\rule{\qpJTcRuleLen}{\qpJTcRule}}}
\newcommand{\jfuti}{{\fontsize{\qpJTdSize}{18pt}\selectfont\heitiao\color{subgrey}%
  \renewcommand{\CJKglue}{\hskip 0pt}第二章}}
\newcommand{\jptime}{{\fontsize{\qpJTeSize}{17.7pt}\selectfont（时间：120分钟\quad 分值：150分）}}
\newcommand{\zuhang}[2]{\par\noindent\hangindent=5.7mm\hangafter=1
  {\fontsize{\qpJTfSize}{17.7pt}\selectfont\heitiao #1}#2\par}
\newcommand{\ti}[2]{\par\noindent\hangindent=5.7mm\hangafter=1
  {\fontsize{11.4pt}{17.7pt}\selectfont\heihao{\numboldjian #1}.}\hspace{2.7mm}#2\par}
\newcommand{\fenzhi}[1]{（#1分）}
\newcommand{\optline}[1]{\par\noindent\hangindent=5.7mm\hangafter=0
  {\fontsize{10.5pt}{19pt}\selectfont #1}\par}
\newcommand{\optII}[2]{\makebox[53.05mm][l]{#1}#2}
\newcommand{\optIV}[4]{\makebox[21.5mm][l]{#1}\makebox[21.5mm][l]{#2}\makebox[21.5mm][l]{#3}#4}
\newcommand{\kw}{\hfill（\hspace{3.6mm}）\hspace*{5.4mm}}
\newcommand{\qpart}[1]{\par\noindent\hangindent=5.7mm\hangafter=0
  \hspace*{5.7mm}{\fontsize{10.5pt}{17.7pt}\selectfont #1}\par}

\newdimen\hA \hA=3.03mm
\newdimen\hB \hB=12.12mm
\newdimen\hC \hC=1.89mm
\newdimen\hD \hD=5.95mm
\newdimen\hE \hE=8.04mm
\newdimen\hF \hF=5.70mm
\newdimen\hG \hG=6.35mm
\newcommand{\jphead}{\hbox to\jpcolw{\sishou\hfill}\nointerlineskip
  \vskip\hA\jpright{\jtianxie}%
  \vskip\hB\jpcen{\jjuanming}%
  \vskip\hC\jpcen{\jjunrule}%
  \vskip\hD\jpcen{\jfuti}%
  \vskip\hE\jpcen{\jptime}%
  \vskip\hF
  \zuhang{一、选择题}{：本题共8小题，每小题5分，共40分．\ 在每小题给出的四个选项中，只有一项是符合题目要求的．}%
  \vgt{\hG}}

% ------------------------------------------------------------
% 件型层 D｜YJ-02 文字行页脚（照抄样张：奇页右置卷名·测评卷，偶页左置羿郭工作室品牌）
% ------------------------------------------------------------
\newcommand{\jpnum}[1]{{\fontsize{\qpYJbNumSize}{\qpYJbNumSize}\selectfont\numbold #1}}
\newcommand{\jptxtsize}{7.5pt}
\newcommand{\jptxt}{\fontsize{\jptxtsize}{\jptxtsize}\selectfont}
\newcommand{\juanname}{单元素养测评卷（二）}
\newcommand{\pqt}{\ifnum\value{page}<10 0\fi\thepage}
\setlength{\headwidth}{\textwidth}
\fancyfoot[RO]{\jptxt\juanname\hspace{3.95mm}{\heitiao 测评卷}\hspace{3.88mm}卷\hspace{2.25mm}\jpnum{\pqt}}
\fancyfoot[LE]{\jptxt\jpnum{\pqt}\hspace{1.9mm}卷\hspace{3.95mm}{\heitiao 羿郭工作室}%
  \hspace{3.0mm}高中数学\hspace{2.75mm}选择性必修第一册\hspace{2.8mm}RJB}

% ------------------------------------------------------------
% 件型层 E′｜三栏排布器（S5-W4 阶段二：零高固定栏 → multicols 三栏回流）
%   观感三参数锁定：栏宽 \jpcolw=120mm（＝multicol 列宽）、栏距 \jpgap=11mm（＝\columnsep，
%   同旧手拼 \hspace{\jpgap}）、无栏线（旧手拼三栏本无线）。全幅 382＝3×120＋2×11。
%   退役（案B）：\jpcol 零高申报／\jpthree 手分栏／\jplead／「\setlength\linewidth 换装补钉」
%   ／卷末附卷页／\anskey 补偿发射／速查表 \datu+\dabiao——答案块随题入流，
%   锚点由 ansblockbegin 恒发（sty:556，双档 M3-ANSKEY 恒等由源结构保证）。
%   [D] 门「零高固定栏承件族」判定随零高结构消失自然失配，卷件退出承件族走 strict 严档。
% ------------------------------------------------------------
\setlength{\columnsep}{\jpgap}
\setlength{\columnseprule}{0pt}
\newenvironment{juancols}
  {\begin{multicols}{3}\fontsize{10.5pt}{17.7pt}\selectfont}%
  {\end{multicols}}

\begin{document}
% ============ S5-W4 阶段二：题-答-解析同流（三栏回流制），案B 手分页/附卷退役 ============
\begin{juancols}
\begin{document}
\jpthree{\jphead
  \ti{1}{已知点 \(A(-1,4)\)，\(B(7,-2)\)，则线段 \(AB\) 的中点坐标与 \(|AB|\) 的值分别为\kw}
  \optline{\optII{A．\((3,1)\)，\(10\)；}{B．\((6,2)\)，\(10\)；}}
  \optline{\optII{C．\((3,1)\)，\(100\)；}{D．\((4,2)\)，\(10\)．}}
  \begin{ansblock}[2章-测-1]
  % ans:2章-测-1
  \ansitem{1}{A}
  \ansnote{解析}{中点 \(M=\left(\dfrac{-1+7}{2},\dfrac{4-2}{2}\right)=(3,1)\)；\(|AB|=\sqrt{8^{2}+(-6)^{2}}=10\)。B 系中点漏除以 \(2\)；C 系忘开方；D 系中点算误。故选 A。}
  \end{ansblock}
  \ti{2}{下列说法中正确的是\kw}
  \optline{A．若两条直线斜率相等，则它们互相平行；}
  \optline{B．若 \(l_{1}\parallel l_{2}\)，则 \(k_{l_{1}}=k_{l_{2}}\)；}
  \optline{C．若两条直线中有一条直线的斜率不存在，另一条直线的斜率存在，则这两条直线相交；}
  \optline{D．若两条直线的斜率都不存在，则它们相互平行．}
  \begin{ansblock}[2章-测-2]
  % ans:2章-测-2
  \ansitem{2}{C}
  \end{ansblock}
  \ti{3}{已知点 \(A(1,4)\)，\(B(5,2)\)，点 \(P\) 为直线 \(l\colon x+y=0\) 上的动点，则 \(|PA|+|PB|\) 的最小值为\kw}
  \optline{\optIV{A．\(3\sqrt{10}\)；}{B．\(2\sqrt{5}\)；}{C．\(\sqrt{34}\)；}{D．\(3\sqrt{5}\)．}}
  \begin{ansblock}[2章-测-3]
  % ans:2章-测-3
  \ansitem{3}{A}
  \ansnote{解析}{\(A\)、\(B\) 代入 \(x+y=0\) 左侧均大于 \(0\)，同侧。设 \(A\) 关于 \(l\) 的对称点 \(A'(-4,-1)\)，则 \(|PA|+|PB|=|PA'|+|PB|\geq|A'B|\)，取等点 \(P_{0}\left(-\dfrac{1}{4},\dfrac{1}{4}\right)\in l\)。最小值 \(=|A'B|=\sqrt{9^{2}+3^{2}}=3\sqrt{10}\)。故选 A。}
  \end{ansblock}
  \ti{4}{与 \(y\) 轴相交于 \(A(0,2)\)、\(B(0,6)\) 两点，且半径等于 \(2\sqrt{2}\) 的圆的方程是\kw}
  \optline{\optII{A．\((x+2)^{2}+(y-4)^{2}=8\)；}{B．\((x-2)^{2}+(y-4)^{2}=8\)；}}
  \optline{C．\((x+2)^{2}+(y-4)^{2}=8\) 或 \((x-2)^{2}+(y-4)^{2}=8\)；}
  \optline{D．\((x-2)^{2}+(y+4)^{2}=8\)．}
  \begin{ansblock}[2章-测-4]
  % ans:2章-测-4
  \ansitem{4}{C（两圆并取）}
  \end{ansblock}
  \ti{5}{已知两点 \(A(3,0)\)，\(B(0,3)\)，动点 \(P\) 满足 \(|PA|^{2}+|PB|^{2}=10\)，则 \(P\) 的轨迹方程为\kw}
  \optline{\optII{A．\(x^{2}+y^{2}-3x-3y+4=0\)；}{B．\(x^{2}+y^{2}-3x-3y-4=0\)；}}
  \optline{\optII{C．\(x^{2}+y^{2}+3x+3y+4=0\)；}{D．\(x^{2}+y^{2}-3x-3y+10=0\)．}}
  \begin{ansblock}[2章-测-5]
  % ans:2章-测-5
  \ansitem{5}{A}
  \ansnote{解析}{设 \(P(x,y)\)：\((x-3)^{2}+y^{2}+x^{2}+(y-3)^{2}=10\)，化简得 \(x^{2}+y^{2}-3x-3y+4=0\)，即 \(\left(x-\dfrac{3}{2}\right)^{2}+\left(y-\dfrac{3}{2}\right)^{2}=\dfrac{1}{2}\)。全过程同解变形，且 \(A\)、\(B\) 均不满足条件，纯粹性与完备性同时成立。故选 A。}
  \end{ansblock}
  \ti{6}{到两定点 \(F_{1}(-3,0)\)，\(F_{2}(3,0)\) 的距离之差的绝对值等于 \(6\) 的点 \(M\) 的轨迹为\kw}
  \optline{\optIV{A．椭圆；}{B．两条射线；}{C．双曲线；}{D．线段．}}
  \begin{ansblock}[2章-测-6]
  % ans:2章-测-6
  \ansitem{6}{B（两条射线）}
  \end{ansblock}
  \ti{7}{若抛物线 \(y^{2}=\dfrac{4}{m}x\) 的焦点与椭圆 \(\dfrac{x^{2}}{7}+\dfrac{y^{2}}{3}=1\) 的左焦点重合，则 \(m\) 的值为\kw}
  \optline{\optIV{A．\(-\dfrac{1}{2}\)；}{B．\(\dfrac{1}{2}\)；}{C．\(-2\)；}{D．\(2\)．}}
  \begin{ansblock}[2章-测-7]
  % ans:2章-测-7
  \ansitem{7}{A（\(m=-\dfrac{1}{2}\)）}
  \end{ansblock}
  \ti{8}{已知 \(F_{1}\)，\(F_{2}\) 分别是双曲线 \(C\colon \dfrac{x^{2}}{3}-\dfrac{y^{2}}{9}=1\) 的左、右焦点，过 \(F_{2}\) 的直线与双曲线 \(C\) 的右支交于 \(A\)，\(B\) 两点，\(\triangle AF_{1}F_{2}\) 和 \(\triangle BF_{1}F_{2}\) 的内心分别为 \(M\)，\(N\)，则 \(|MN|\) 的取值范围是\kw}
  \optline{\optIV{A．\([2\sqrt{3},4)\)；}{B．\([2\sqrt{3},3)\)；}{C．\([2\sqrt{3},+\infty)\)；}{D．\((2\sqrt{3},3)\)．}}
  \begin{ansblock}[2章-测-8]
  % ans:2章-测-8
  \ansitem{8}{A（\([2\sqrt{3},4)\)）}
  \end{ansblock}
  \zuhang{二、选择题}{：本题共3小题，每小题6分，共18分．\ 在每小题给出的选项中，有多项符合题目要求．}%
  \vgt{\hG}
  \ti{9}{\fenzhi{6}点 \(M(1,1)\) 到抛物线 \(y=ax^{2}\) 的准线的距离为 \(2\)，则 \(a\) 的值可以为\kw}
  \optline{\optIV{A．\(\dfrac{1}{4}\)；}{B．\(-\dfrac{1}{12}\)；}{C．\(\dfrac{1}{12}\)；}{D．\(-\dfrac{1}{4}\)．}}
  \begin{ansblock}[2章-测-9]
  % ans:2章-测-9
  \ansitem{9}{AB}
  \end{ansblock}
  \ti{10}{\fenzhi{6}已知 \(F_{1}\)、\(F_{2}\) 分别为双曲线 \(\dfrac{x^{2}}{a^{2}}-\dfrac{y^{2}}{b^{2}}=1\)（\(a>0\)，\(b>0\)）的左、右焦点，且 \(|F_{1}F_{2}|=\dfrac{2b^{2}}{a}\)，点 \(P\) 为双曲线右支一点，\(I\) 为 \(\triangle PF_{1}F_{2}\) 的内心，若 \(S_{\triangle IPF_{1}}=S_{\triangle IPF_{2}}+\lambda S_{\triangle IF_{1}F_{2}}\) 成立，则下列结论正确的有\kw}
  \optline{A．当 \(PF_{2}\perp x\) 轴时，\(\angle PF_{1}F_{2}=30^{\circ}\)；}
  \optline{B．离心率 \(e=\dfrac{1+\sqrt{5}}{2}\)；}
  \optline{C．\(\lambda=\dfrac{\sqrt{5}-1}{2}\)；}
  \optline{D．点 \(I\) 的横坐标为定值 \(a\)．}
  \begin{ansblock}[2章-测-10]
  % ans:2章-测-10
  \ansitem{10}{BCD}
  \end{ansblock}
  \ti{11}{\fenzhi{6}已知线段 \(BC\) 的长度为 \(4\)，线段 \(AB\) 的长度为 \(m\)，点 \(D\)，\(G\) 满足 \(\overrightarrow{AD}=\overrightarrow{DC}\)，\(\overrightarrow{DG}\cdot\overrightarrow{AC}=0\)，且 \(G\) 点在直线 \(AB\) 上，若以 \(BC\) 所在直线为 \(x\) 轴，\(BC\) 的中垂线为 \(y\) 轴建立平面直角坐标系，则\kw}
  \optline{A．当 \(m=4\) 时，点 \(G\) 的轨迹为圆；}
  \optline{B．当 \(6\leq m\leq 8\) 时，点 \(G\) 的轨迹为椭圆，且椭圆的离心率取值范围为 \(\left[\dfrac{1}{2},\dfrac{2}{3}\right]\)；}
  \optline{C．当 \(m=2\) 时，点 \(G\) 的轨迹为双曲线，且该双曲线的渐近线方程为 \(y=\pm\sqrt{3}x\)；}
  \optline{D．当 \(m=5\) 时，\(\triangle BCG\) 面积的最大值为 \(3\)．}
  \begin{ansblock}[2章-测-11]
  % ans:2章-测-11
  \ansitem{11}{BCD}
  \end{ansblock}
  \vgt{\hG}
  \zuhang{三、填空题}{：本题共3小题，每小题5分，共15分．}%
  \vgt{\hG}
  \ti{12}{\fenzhi{5}抛物线 \(x^{2}=-4y\) 的准线方程为\kongbai}
  \begin{ansblock}[2章-测-12]
  % ans:2章-测-12
  \ansitem{12}{\(y=1\)}
  \end{ansblock}
  \ti{13}{\fenzhi{5}过抛物线 \(y^{2}=4x\) 的焦点 \(F\) 作垂直于 \(x\) 轴的直线，交抛物线于 \(A\)，\(B\) 两点，则以 \(F\) 为圆心，\(AB\) 为直径的圆的方程是\kongbai}
  \begin{ansblock}[2章-测-13]
  % ans:2章-测-13
  \ansitem{13}{\((x-1)^{2}+y^{2}=4\)}
  \end{ansblock}
  \ti{14}{\fenzhi{5}双曲线 \(C\colon \dfrac{x^{2}}{a^{2}}-\dfrac{y^{2}}{b^{2}}=1\)（\(a>0\)，\(b>0\)）的左、右焦点分别为 \(F_{1}\)，\(F_{2}\)，右支上有一点 \(M\)，满足 \(\angle F_{1}MF_{2}=90^{\circ}\)，\(\triangle F_{1}MF_{2}\) 的内切圆与 \(y\) 轴相切，则双曲线 \(C\) 的离心率为\kongbai}
  \begin{ansblock}[2章-测-14]
  % ans:2章-测-14
  \ansitem{14}{\(1+\sqrt{3}\)}
  \end{ansblock}
  \vgt{\hG}
  \zuhang{四、解答题}{：本题共5小题，共77分．\ 解答应写出文字说明、证明过程或演算步骤．}%
  \vgt{\hG}
  \ti{15}{\fenzhi{13}已知 \(P(1,2)\) 在抛物线 \(C\colon y^{2}=2px\) 上．}
  \qpart{（1）求抛物线 \(C\) 的方程；}
  \qpart{（2）\(A\)，\(B\) 是抛物线 \(C\) 上的两个动点，如果直线 \(PA\) 的斜率与直线 \(PB\) 的斜率之和为 \(2\)，证明：直线 \(AB\) 过定点．}
  \begin{ansblock}[2章-测-15]
  % ans:2章-测-15
  \ansitem{15}{(1)\(y^{2}=4x\)；(2)直线 \(AB\) 恒过定点 \((-1,0)\)}
  \ansnote{解析}{(1) \(4=2p\)，\(p=2\)，\(y^{2}=4x\)。(2) 设 \(AB\)：\(x=my+t\)，与 \(y^{2}=4x\) 联立得 \(y^{2}-4my-4t=0\)，\(y_{1}+y_{2}=4m\)，\(y_{1}y_{2}=-4t\)（\(\Delta>0\) 须 \(m^{2}+t>0\)）。\(k_{PA}=\dfrac{4}{y_{1}+2}\)，\(k_{PB}=\dfrac{4}{y_{2}+2}\)，由 \(k_{PA}+k_{PB}=2\) 得 \(y_{1}y_{2}=4\)，故 \(-4t=4\)，\(t=-1\)，直线 \(AB\) 恒过定点 \((-1,0)\)。}
  \end{ansblock}
  \ti{16}{\fenzhi{14}已知椭圆 \(C\)：\(\dfrac{x^{2}}{a^{2}}+\dfrac{y^{2}}{3}=1\)（\(a>\sqrt{3}\)）的离心率为 \(\dfrac{\sqrt{2}}{2}\)．}
  \qpart{（1）求椭圆 \(C\) 的方程；}
  \qpart{（2）若直线 \(l\) 经过 \(C\) 的左焦点 \(F_{1}\) 且与 \(C\) 相交于 \(B\)、\(D\) 两点，以线段 \(BD\) 为直径的圆经过椭圆 \(C\) 的右焦点 \(F_{2}\)，求 \(l\) 的方程．}
  \begin{ansblock}[2章-测-16]
  % ans:2章-测-16
  \ansitem{16}{(1)\(\dfrac{x^{2}}{6}+\dfrac{y^{2}}{3}=1\)；(2)\(x-\sqrt{7}y+\sqrt{3}=0\) 或 \(x+\sqrt{7}y+\sqrt{3}=0\)}
  \ansnote{解析}{(1) \(b=\sqrt{3}\)，\(\dfrac{c}{a}=\dfrac{\sqrt{2}}{2}\)，\(a^{2}=b^{2}+c^{2}\) 解得 \(a=\sqrt{6}\)。(2) \(F_{1}(-\sqrt{3},0)\)、\(F_{2}(\sqrt{3},0)\)，设 \(l\)：\(x=my-\sqrt{3}\)，代入得 \((m^{2}+2)y^{2}-2\sqrt{3}my-3=0\)（\(\Delta>0\) 恒成立），\(\overrightarrow{F_{2}B}\cdot\overrightarrow{F_{2}D}=0\) 即 \(x_{1}x_{2}-\sqrt{3}(x_{1}+x_{2})+y_{1}y_{2}+3=0\)，代入韦达整理得 \(m^{2}=7\)，\(m=\pm\sqrt{7}\)，\(l\)：\(x-\sqrt{7}y+\sqrt{3}=0\) 或 \(x+\sqrt{7}y+\sqrt{3}=0\)。}
  \end{ansblock}
  \ti{17}{\fenzhi{15}椭圆 \(\dfrac{x^{2}}{a^{2}}+\dfrac{y^{2}}{b^{2}}=1\)（\(a>b>0\)）离心率为 \(\dfrac{\sqrt{3}}{3}\)，直线 \(x-2y+b=0\) 与椭圆交于 \(P\)，\(Q\) 两点，且 \(PQ\) 中点为 \(E\)，\(O\) 为原点，求直线 \(OE\) 的斜率．}
  \begin{ansblock}[2章-测-17]
  % ans:2章-测-17
  \ansitem{17}{\(-\dfrac{4}{3}\)}
  \ansnote{解析}{由 \(e=\dfrac{\sqrt{3}}{3}\) 得 \(\dfrac{b^{2}}{a^{2}}=\dfrac{2}{3}\)。设 \(P(x_{1},y_{1})\)、\(Q(x_{2},y_{2})\)，\(k_{PQ}=\dfrac{1}{2}\)，两式作差得 \(\dfrac{(y_{1}-y_{2})(y_{1}+y_{2})}{(x_{1}-x_{2})(x_{1}+x_{2})}=-\dfrac{b^{2}}{a^{2}}=-\dfrac{2}{3}\)，即 \(k_{PQ}\cdot k_{OE}=-\dfrac{2}{3}\)，故 \(k_{OE}=-\dfrac{4}{3}\)。回代验 \(\Delta\)：联立化简得 \(11x^{2}+6bx-9b^{2}=0\)，\(\Delta=432b^{2}>0\)，弦恒存在。}
  \end{ansblock}
  \ti{18}{\fenzhi{16}在平面直角坐标系 \(xOy\) 中，已知点 \(M\) 是双曲线 \(\dfrac{x^{2}}{a^{2}}-\dfrac{y^{2}}{b^{2}}=1\)（\(a>0\)，\(b>0\)）上的异于顶点的任意一点，过点 \(M\) 作双曲线的切线 \(l\)，若 \(k_{OM}\cdot k_{l}=\dfrac{1}{3}\)，求双曲线离心率 \(e\) 的值．}
  \begin{ansblock}[2章-测-18]
  % ans:2章-测-18
  \ansitem{18}{\(\dfrac{2\sqrt{3}}{3}\)}
  \ansnote{解析}{当 \(y>0\) 时 \(y'=\dfrac{b^{2}x}{a^{2}y}\)，得切线方程 \(\dfrac{x_{0}}{a^{2}}x-\dfrac{y_{0}}{b^{2}}y=1\)（\(y\leq 0\) 时同理成立）。设 \(M(m,n)\)：\(k_{l}=\dfrac{b^{2}m}{a^{2}n}\)，\(k_{OM}=\dfrac{n}{m}\)，由 \(k_{OM}\cdot k_{l}=\dfrac{1}{3}\) 得 \(\dfrac{b^{2}}{a^{2}}=\dfrac{1}{3}\)，故 \(e=\sqrt{1+\dfrac{1}{3}}=\dfrac{2\sqrt{3}}{3}\)。}
  \end{ansblock}
  \ti{19}{\fenzhi{19}已知双曲线 \(\dfrac{x^{2}}{a^{2}}-\dfrac{y^{2}}{3a^{2}}=1\) 与椭圆 \(\dfrac{x^{2}}{a^{2}}+\dfrac{y^{2}}{b^{2}}=16\)（\(a>b>0\)）共焦点，且双曲线与直线 \(y=4bx-3\) 相切，求 \(a\) 的值．}
  \begin{ansblock}[2章-测-19]
  % ans:2章-测-19
  \ansitem{19}{\(a=1\)}
  \ansnote{解析}{共焦点：\(3a^{2}+a^{2}=16(a^{2}-b^{2})\)，即 \(3a^{2}=4b^{2}\)。联立 \(y=4bx-3\) 化简得 \((3-16b^{2})x^{2}+24bx-9-3a^{2}=0\)，相切即 \(\Delta=0\)：\(48b^{2}+(3-16b^{2})(3+a^{2})=0\)，代入 \(3a^{2}=4b^{2}\) 得 \(4a^{4}-a^{2}-3=0\)，\(a^{2}=1\)，\(a=1\)（负值舍去）。}
  \end{ansblock}
\end{juancols}
\end{document}
